% ============================================================
% Section 75: 一维双原子链的紧束缚近似
% ============================================================
\section{固体理论：一维双原子链的紧束缚近似}
\label{sec:1d-diatomic-tight-binding}

\begin{modelbox}[题目：一维双原子链的紧束缚能带]
由相同原子组成的一维原子链，每个原胞中有两个原子，原胞长度为 $a$，原胞内两个原子相对距离为 $b$。
\begin{enumerate}
    \item 根据紧束缚近似，只计入近邻相互作用，写出原子 $s$ 态相对应的晶体波函数形式；
    \item 求出相应能带的 $E(k)$ 函数。
\end{enumerate}
\end{modelbox}

\begin{tipbox}[考点快速分析]
\begin{itemize}
    \item \textbf{复式格子}：两套子晶格，需用两个布洛赫和线性组合
    \item \textbf{久期方程}：$2\times2$ 矩阵，非对角元为子晶格间跃迁积分
    \item \textbf{两支能带}：成键带（低能）和反键带（高能），在布里渊区边界可能打开能隙
    \item \textbf{对称性极限}：$b=a/2$ 时两支能带在边界简并（退化为单原子链）
\end{itemize}
\end{tipbox}

\begin{derivationbox}[标准解题过程]
\textbf{Step 1: 晶体波函数的形式}

这是一维复式格子，包含两套等价的子晶格：
\begin{itemize}
    \item A 子晶格：原子位于 $X_n^{(A)}=na$
    \item B 子晶格：原子位于 $X_n^{(B)}=na+b$
\end{itemize}

对 $s$ 态原子轨道 $\varphi_s(x)$，两套子晶格的布洛赫和分别为
\begin{align}
\Psi_{k,A}(x) &= \frac{1}{\sqrt{N}}\sum_{n}e^{ikna}\varphi_s(x-na),\\
\Psi_{k,B}(x) &= \frac{1}{\sqrt{N}}\sum_{n}e^{ik(na+b)}\varphi_s(x-na-b).
\end{align}

晶体的电子波函数是两者的线性组合：
\begin{equation}
\boxed{\Psi_k(x) = A\,\Psi_{k,A}(x) + B\,\Psi_{k,B}(x),}
\end{equation}
其中 $A,B$ 为待定系数，由薛定谔方程确定。

\textbf{Step 2: 久期方程的建立}

将 $\Psi_k$ 代入 $H\Psi_k=E\Psi_k$，分别左乘 $\Psi_{k,A}^*$ 和 $\Psi_{k,B}^*$ 并积分，得到关于 $A,B$ 的齐次线性方程组。非零解要求久期行列式为零：
\begin{equation}
\begin{vmatrix}
H_{11}-ES_{11} & H_{12}-ES_{12} \\[4pt]
H_{21}-ES_{21} & H_{22}-ES_{22}
\end{vmatrix}=0,
\end{equation}
其中 $H_{ij}=\langle\Psi_{k,i}|H|\Psi_{k,j}\rangle$，$S_{ij}=\langle\Psi_{k,i}|\Psi_{k,j}\rangle$。

\textbf{紧束缚近似下的简化}：
\begin{itemize}
    \item 同一原子轨道归一化：$S_{11}=S_{22}=1$
    \item 不同格点轨道近似正交：$S_{12}=S_{21}=0$
    \item 对角元（原地能量，含最近邻平均势场）：$H_{11}=H_{22}=E_0$
\end{itemize}

非对角元 $H_{12}$ 描述 A、B 子晶格间的跃迁。每个 A 原子有两个最近邻 B 原子：
\begin{itemize}
    \item 同原胞内 B：距离 $b$，相位因子 $e^{ikb}$
    \item 左侧原胞 B：距离 $a-b$，相位因子 $e^{-ik(a-b)}$
\end{itemize}
设对应的跃迁积分分别为 $t_1$ 和 $t_2$（通常 $t_1,t_2<0$），则
\begin{equation}
H_{12}=t_1 e^{ikb}+t_2 e^{-ik(a-b)},\qquad H_{21}=H_{12}^*.
\end{equation}

\textbf{Step 3: 能带求解}

久期方程化为
\begin{equation}
\begin{vmatrix}
E_0-E & H_{12} \\[4pt]
H_{12}^* & E_0-E
\end{vmatrix}=0
\quad\Longrightarrow\quad
(E_0-E)^2-|H_{12}|^2=0.
\end{equation}

解得两支能带：
\begin{equation}
E(k)=E_0\pm|H_{12}|=E_0\pm\bigl|t_1 e^{ikb}+t_2 e^{-ik(a-b)}\bigr|.
\end{equation}

提取公共相位并化简（设 $t_1,t_2$ 为实数）：
\begin{align}
|H_{12}|^2 &= \bigl(t_1 e^{ikb}+t_2 e^{-ik(a-b)}\bigr)\bigl(t_1 e^{-ikb}+t_2 e^{ik(a-b)}\bigr) \notag\\
&= t_1^2+t_2^2+t_1t_2\bigl(e^{ika}+e^{-ika}\bigr) \notag\\
&= t_1^2+t_2^2+2t_1t_2\cos(ka).
\end{align}

因此一般结果为
\begin{equation}
\boxed{E(k)=E_0\pm\sqrt{t_1^{\,2}+t_2^{\,2}+2t_1t_2\cos(ka)}.}
\end{equation}

\textbf{重要极限}：
\begin{itemize}
    \item \textbf{$t_1=t_2=t$}（原胞内与跨原胞跃迁相等，如 $b=a/2$ 时由对称性保证）：
    \begin{equation}
    E(k)=E_0\pm\sqrt{2t^2+2t^2\cos(ka)}=E_0\pm 2|t|\Bigl|\cos\frac{ka}{2}\Bigr|.
    \end{equation}
    在布里渊区边界 $k=\pi/a$ 处，$\cos(ka/2)=0$，两支能带简并（退化为单原子链）。
    \item \textbf{$t_2=0$}（只考虑原胞内跃迁）：
    \begin{equation}
    E(k)=E_0\pm|t_1|,\quad\text{平带，无色散。}
    \end{equation}
    这说明跨原胞的跃迁是形成能带色散的必不可少的因素。
\end{itemize}
\end{derivationbox}

\begin{formulabox}[答案汇总]
\begin{center}
\begin{tabular}{ll}
\toprule
\textbf{问题} & \textbf{答案} \\
\midrule
(1) 波函数 & $\Psi_k=A\Psi_{k,A}+B\Psi_{k,B}$（两套子晶格布洛赫和线性组合） \\
(2) 能带 & $E(k)=E_0\pm\sqrt{t_1^{\,2}+t_2^{\,2}+2t_1t_2\cos(ka)}$ \\
\quad $t_1=t_2=t$ 时 & $E(k)=E_0\pm 2|t|\,|\cos(ka/2)|$ \\
\bottomrule
\end{tabular}
\end{center}
\end{formulabox}

\begin{insightbox}[物理图像]
\begin{itemize}
    \item 一维复式格子的紧束缚能带分裂为\textbf{两支}，对应原胞内两个原子的成键组合（$-$ 号，低能）和反键组合（$+$ 号，高能）。
    \item 能带宽度由跃迁积分 $|t_1|,|t_2|$ 决定。若 $|t_1|\gg|t_2|$（原子在原胞内靠得很近），能带几乎为平带，电子被``囚禁''在原胞内。
    \item 布里渊区边界 $k=\pi/a$ 处的能隙大小为 $2|t_1-t_2|$（由 $E_+(\pi/a)-E_-(\pi/a)=2||t_1|-|t_2||$ 给出）。只有当 $t_1=t_2$ 时能隙闭合——这正是 $b=a/2$ 时对称性提高的结果。
    \item 本题与 sec74 的 $\delta$ 函数势模型形成对照：二者都研究一维双原子链，但 sec74 用近自由电子近似（弱周期势），本题用紧束缚近似（强原子束缚）。在 $b=a/2$ 的极限下，两者都预言了从绝缘体到金属的转变。
\end{itemize}
\end{insightbox}
